\subsection{Bosonic relative index and Higgs selection}

\begin{theorem}[Minimal determinant classes from the primitive charge generator]
\label{thm:determinant-class-higgs-selection}
\label{thm:determinant-clutching}
Let
\[
L_2:=\det E_2,
\qquad
L_3:=\det E_3
\]
over the compactified oriented normal sphere $S^2=D_N\cup_{S^1}D_S$. Then the primitive charge
generator
\[
X=-2P_3+3P_2
\]
together with center neutrality and seam-evenness forces the clutching map
\[
g_2(\varphi)=e^{i\varphi}
\]
for $L_2$, while center neutrality forces
\[
g_3(\varphi)=1
\]
for $L_3$. Therefore
\[
c_1(L_2)=\deg g_2=1,
\qquad
c_1(L_3)=\deg g_3=0.
\]
This pair is the unique minimal nonnegative admissible determinant class.
\end{theorem}

\begin{proof}
Complex line bundles on
\[
S^2=D_N\cup_{S^1}D_S
\]
are classified by their clutching maps along the equator:
\[
\pi_1(U(1))\cong \mathbb Z.
\]
Hence the first Chern class is exactly the winding degree of the transition function.

In the weak rank-two sector the primitive charge generator
\[
X=-2P_3+3P_2
\]
assigns the unique positive determinant winding to the seam-even weak block. The spin lift on the
oriented normal sphere acts there by the half-angle representation, so on the determinant line the
transition factor is
\[
g_2(\varphi)=e^{i\varphi},
\]
hence
\[
\deg g_2=1,
\qquad
c_1(L_2)=1.
\]

In the color sector admissibility imposes center neutrality. Therefore the determinant transition
function of the color block is null-homotopic and may be chosen constant:
\[
g_3(\varphi)=1.
\]
Consequently
\[
\deg g_3=0,
\qquad
c_1(L_3)=0.
\]

The full primitive seam twist is therefore exhausted in the weak determinant class. Any nonnegative
pair strictly smaller than $(1,0)$ would force vanishing weak determinant winding and hence remove
the seam-even weak line required by the primitive charge generator. Any larger nonnegative pair
would insert extra clutching not carried by the primitive generator and therefore violate
minimality. Thus the unique minimal nonnegative admissible determinant class is
\[
(c_1(L_2),c_1(L_3))=(1,0).
\]
\end{proof}

\begin{remark}[Energy interpretation on the retained branch]
The older asymptotic bosonic functional
\[
\mathcal E_{\mathrm{asym}}(\Phi_2,\Phi_3)
:=
\tau_2\bigl|\mathrm{wind}(\det\Phi_2)\bigr|
+\tau_3\bigl|\mathrm{wind}(\det\Phi_3)\bigr|
+m_3^2\dim\ker\mathcal B_{E_3}^+,
\]
with $\tau_3>\tau_2>0$, remains a useful physical intuition for why the retained branch is energetically preferred. In the present rewrite, however, the structural weight is carried by the determinant-class theorem above rather than by the cost functional itself.
\end{remark}

\begin{theorem}[Bosonic index and unique Higgs doublet from the minimal determinant class]
\label{thm:callias-higgs-selection}
\label{thm:compact-higgs-unique}
\label{thm:compact-determinant-higgs-index}
Compactify the oriented two-dimensional normal slice $N_\Seam$ to $S^2$. For $a\in\{2,3\}$
define
\[
\mathcal B_{E_a}:=\sigma^\mu\nabla^{(a)}_\mu+\Phi_a
\]
on the compactified normal sphere. Under the determinant-class data of
\Cref{thm:determinant-class-higgs-selection}, one has
\[
L_2\cong\mathcal O(1),
\qquad
L_3\cong\mathcal O,
\]
and hence
\[
H^0(S^2,\mathcal O(1))\cong\mathbb C^2,
\qquad
H^1(S^2,\mathcal O(1))=0,
\]
and the compact bosonic index satisfies
\[
\Ind(\mathcal B_{E_2}^+)=1,
\qquad
\Ind(\mathcal B_{E_3}^+)=0.
\]
Consequently the closed branch contains exactly one complex weak doublet and no light color triplet,
and the unique light seam-even bosonic zero mode lies in
\[
\Phi\in\Gamma(E_2)\cong (1,2)_{1/2},
\qquad
\taudbl^{\ast}\Phi=+\Phi.
\]
\end{theorem}

\begin{proof}
By \Cref{thm:determinant-class-higgs-selection}, the weak determinant line is the degree-one line
bundle
\[
L_2\cong\mathcal O(1)
\]
on
\[
S^2\cong\mathbb CP^1,
\]
whereas the color determinant class is trivial and
\[
L_3\cong\mathcal O.
\]
For $\mathbb CP^1$, Birkhoff--Grothendieck identifies every holomorphic line bundle by its degree,
and Riemann--Roch gives
\[
\chi(\mathcal O(1))=1+1=2.
\]
Serre duality yields
\[
H^1(S^2,\mathcal O(1))
\cong
H^0(S^2,\mathcal O(-3))^\vee
=0,
\]
so
\[
H^0(S^2,\mathcal O(1))\cong\mathbb C^2,
\qquad
H^1(S^2,\mathcal O(1))=0.
\]

The resulting two-dimensional zero-mode space is exactly one complex weak doublet, so the compact
bosonic index records it as
\[
\Ind(\mathcal B_{E_2}^+)=1.
\]

In the color sector one has $L_3\cong\mathcal O$, so no positive determinant degree is available.
Center neutrality forbids an unpaired seam-even color contribution, hence the positive and negative
triplet modes pair with zero net index:
\[
\Ind(\mathcal B_{E_3}^+)=0.
\]
Therefore the closed branch contains exactly one seam-even Higgs doublet and no light color
triplet, and the unique light bosonic zero mode lies in the weak block.
\end{proof}

\begin{corollary}[Bosonic rank is two]
\label{cor:bosonic-rank-two}
On the closed branch the positive carrier block has rank
\[
\dim E_+=2.
\]
Equivalently, the unique seam-even light bosonic zero mode fills one complex weak doublet and no
larger positive carrier block is retained.
\end{corollary}

\begin{proof}
By \Cref{thm:callias-higgs-selection}, the light seam-even bosonic zero mode lies in the weak block
\[
\Phi\in\Gamma(E_2)\cong (1,2)_{1/2},
\]
and the compact index identifies its zero-mode space with
\[
H^0(S^2,\mathcal O(1))\cong\mathbb C^2.
\]
Hence the retained positive carrier block is exactly two-dimensional.
\end{proof}

This is the sharp form of the Higgs claim used here. The issue is not whether a Higgs doublet can
be written down, but whether a compact bosonic index can be shown to select it uniquely.
\begin{remark}[Callias shadow]
The older Callias formulation is retained only as the noncompact shadow of the same compact index
theorem and is discussed separately in the appendical analytic remarks. It no longer carries the
main proof burden.
\end{remark}

Once that uniqueness statement is fixed, the renormalizable carrier-compatible Yukawa couplings
are the standard ones,
\[
Q\,u^c\,\Phi,
\qquad
Q\,d^c\,\Phi^\dagger,
\qquad
L\,e^c\,\Phi^\dagger,
\qquad
L\,\nu^c\,\Phi,
\]
so the Higgs statement becomes a genuine structural bridge rather than a notation choice. The
formal carrier discharge from this Yukawa bridge is recorded later in
\Cref{cor:bosonic-rank-two,thm:yukawa-forces-32,cor:yukawa-discharges-kpp}, once the closure-theorem Higgs block has
been stated in its final retained-branch form.

\subsection{Closed \texorpdfstring{$41$}{41} chain after family and Higgs closure}

Once the derived corner theorem and the Higgs uniqueness theorem are in place, the generic
abelian-index identity \Cref{eq:41-chain-main} specializes to
\begin{equation}
\label{eq:41-chain-specialized}
10\,\bone
=
\Oadm\gamma+1
=
48\cdot\frac56+1
=
41.
\end{equation}
The extra $+1$ is precisely the arithmetic contribution of the unique light seam-even Higgs
doublet.

\begin{figure}[t]
\centering
\footnotesize
\begin{tikzpicture}[
    >=Latex,
    chain/.style={draw=black!55, rounded corners, align=center, text width=2.6cm, minimum height=1.0cm, inner sep=4pt},
    eq/.style={font=\Large},
    note/.style={font=\scriptsize\itshape, text=black!55, align=center}
]
\node[chain, fill=green!7, draw=green!50!black] (a) at (0,0) {$10\,\bone$\\\scriptsize abelian index};
\node[eq] at (2.15,0) {$=$};
\node[chain, fill=orange!9, draw=orange!75!black] (b) at (4.3,0) {$\Oadm\gamma+1$\\\scriptsize family occupancy $+$ Higgs};
\node[eq] at (6.45,0) {$=$};
\node[chain, fill=gray!8, draw=black!65] (c) at (8.6,0) {$41$\\\scriptsize closed specialization};
\node[note] at (4.3,-1.2) {This specialization appears only after both $\,\Oadm=48\,$ and $\,N_\Phi=1\,$ are proved.};
\end{tikzpicture}
\caption{The closed $41$ chain. The arithmetic specialization is made only after the family and
Higgs theorems fix the occupancy and the unique light doublet.}
\label{fig:41-chain-strip}
\end{figure}

